\begin{theorem}[E.34 --- Proved: full metabelian zero-abelian families at every five-power level]
\label{E-five-metabelian}
Let $n=5^r$, $r\ge1$, and let
$N=\pi_1(U_{n,\overline{\mathbb Q}})$ be the full
profinite geometric group with the rational unit
tangent $s_O$ at $O=[0:-1:1]$. Put
\[
 M=N/N'',\qquad d_n=\frac{(n^2-1)(n^2+3)}4.
\]
There is a continuous injective family
\[
 \mathbb Z_5^{d_n}\longrightarrow
 \operatorname{Sec}(M\rtimes_{s_O}G_{\mathbb Q})/M
 \tag{E.34.1}
\]
with zero abelian cocycle. Its zero member is
$[s_O]$, and every nonzero member is noncuspidal,
including for the cusp normalizers in this
metabelian quotient itself. In particular
$d_5=168$.
\end{theorem}
\begin{proof}
Put $P=N^{(5)}$, $V=P^{\mathrm{ab}}$,
$D_5=P'/P''$, and $B_2=\bigwedge^2V$.
The same characteristic Haar summand as in
E.31, now in the coefficient-five factor,
gives a continuous arithmetic-equivariant
embedding
$j:D_5\hookrightarrow N'/N''$ which is a
section of the natural comparison. The
surjection $D_5\twoheadrightarrow B_2$ is
the actual degree-two quotient.

E.16 identifies the actual arithmetic action
with a quotient
\[
 \Gamma=P_{\mathrm{ar}}\rtimes C_4,
 \qquad
 P_{\mathrm{ar}}=G_{\mathbb Q(\zeta_5),\{5\}}^{(5)},
 \tag{E.34.2}
\]
where $P_{\mathrm{ar}}$ is free pro-five of
rank three. The complement is chosen by
the profinite coprime splitting theorem;
changing it gives conjugate data. No lift
of that finite complement to $G_{\mathbb Q}$
is required.

We prove surjectivity on cohomology for the
precise comparison
\[
 H^1_{\mathrm{cts}}(\Gamma,D_5)
       \longrightarrow H^1_{\mathrm{cts}}(\Gamma,B_2).
 \tag{E.34.3}
\]
Let $b$ represent a class on the right.
Since four is invertible on these modules,
averaging on $C_4$ changes $b$ by a
coboundary so that $b|_{C_4}=0$.
Its restriction to $P_{\mathrm{ar}}$ is
then exactly invariant for the action
\[
 (a\cdot b)(g)=a\,b(a^{-1}ga),\qquad a\in C_4.
 \tag{E.34.4}
\]
Lift its values on the three free
generators arbitrarily to $D_5$.
Freeness gives a continuous $D_5$-cocycle
$\widetilde b$ on $P_{\mathrm{ar}}$.
The sum
\[
 \widetilde b_{\mathrm{av}}
       =\frac14\sum_{a\in C_4}a\cdot\widetilde b
 \tag{E.34.5}
\]
is meaningful in this abelian pro-five
module. It is a continuous invariant
cocycle, and its projection is exactly
$b|_{P_{\mathrm{ar}}}$ because that
cocycle was already invariant. Setting
its value on $C_4$ equal to zero extends
it to $\Gamma$. This proves (E.34.3)
on actual cocycles, without an inverse
limit of unrelated lift choices.

We compute the rank of the target using
the same arithmetic cohomology inputs
proved in E.16. Explicitly, a $B_2$-cocycle
on $G_{\mathbb Q(\zeta_5),\{5\}}$ has
graph in the pro-five group
$B_2\rtimes\rho(G_{\mathbb Q(\zeta_5),\{5\}})$.
It kills the kernel of the maximal
pro-five quotient. Consequently inflation
identifies, rather than merely embeds,
\[
 H^1(\Gamma,B_2)
       =H^1(G_{\mathbb Q,\{5,\infty\}},B_2).
 \tag{E.34.6}
\]
The finite-module filtration, the
coprime averaging and the finite-level
Mittag--Leffler argument in E.16 give
zero continuous $H^2$ here. Good-Frobenius
weights give $H^0=0$. The global Euler
formula, in its real-place form used
there, therefore gives
\[
 \operatorname{rank}_{\mathbb Z_5}H^1(\Gamma,B_2)
       =\dim_{\mathbb Q_5}(B_2\otimes\mathbb Q_5)^-.
 \tag{E.34.7}
\]
The cusp permutation and compact
intersection-pairing calculation of
(E.33.5) applies to every odd $n$.
Thus the last dimension is $d_n$.

Choose a free lattice of rank $d_n$
in this cohomology group, choose
representative cocycles, and lift
them by (E.34.3). A linear combination
with coefficients in $\mathbb Z_5$
gives a jointly continuous family of
cocycles in $D_5$. Apply $j$ and
multiply by $s_O$. The conjugacy proof
of E.33 applies with five in place
of three: $H^0(G_{\mathbb Q},N^{\rm ab})=0$
forces every conjugator into $N'/N''$,
and projection to $B_2$ then proves
injectivity of the parameter family.

The primitive-loop centralizer proof
in E.33 uses the full completed
group ring and is independent of
the coefficient prime, so it
identifies the cusp normalizers
here as well. Each nonorigin
rational cusp has a nonzero
Abel--Jacobi class killed by
$n=5^r$, by the tangent-line
divisor. Its five-primary Kummer
class is nonzero. Its packet
therefore has no zero-abelian
section. The origin packet has
only $s_O$ at zero abelianization,
and the nonrational cusps have
proper arithmetic stabilizers.
This proves every assertion.
The anchors are \S51.16,
folios733v--736v, scans172--178,
and LM7-E.23(iii).
\end{proof}
